\documentclass[10pt]{article}
\usepackage[a4paper,margin=12mm]{geometry}
\usepackage{fontspec}
\usepackage{graphicx}
\usepackage{enumitem}
\usepackage{multicol}
\usepackage{ragged2e}
\usepackage[most]{tcolorbox}
\usepackage{xcolor}
\usepackage{hyperref}
\setmainfont{Sarabun}
\setsansfont{Sarabun}
\IfFontExistsTF{DejaVu Sans Mono}{\setmonofont{DejaVu Sans Mono}}{\setmonofont{Latin Modern Mono}}
\XeTeXlinebreaklocale "th"
\XeTeXlinebreakskip = 0pt plus 1pt
\setlength{\parindent}{0pt}
\setlength{\parskip}{0.45em}
\tolerance=9999
\emergencystretch=4em
\sloppy
\pagestyle{empty}
\definecolor{panel}{HTML}{FAF7F0}
\definecolor{line}{HTML}{CFC7B6}
\definecolor{muted}{HTML}{5A5A5A}
\definecolor{highlight}{HTML}{8B6F00}
\begin{document}
\section*{Slide 1: Title / Framing}
\noindent\textcolor{muted}{Slide 1 \textbar\ Slide1.md}
\vspace{0.55em}
\begin{center}
\includegraphics[width=\linewidth,height=0.340\textheight,keepaspectratio]{/mnt/c/Users/punpi/OneDrive/KU/president/vision/qna/presentation-narrator/assets/presentation\_slides/Slide1.JPG}
\end{center}
\vspace{0.4em}
\begin{tcolorbox}[breakable,colback=panel,colframe=line,boxrule=0.5pt,arc=2mm,left=3mm,right=3mm,top=2mm,bottom=2mm]
\textbf{Version To Actually Say}
\medskip
\setlength{\columnsep}{1.1em}
\raggedcolumns
\begin{multicols}{2}
\RaggedRight
\setlength{\parskip}{0.7em}
\par วันนี้ผมไม่ได้เริ่มด้วยโครงการ\\ แต่เริ่มด้วย \textbf{\textcolor{highlight}{“วิธีคิด”}} สำหรับระยะถัดไป\\ ของมหาวิทยาลัยเกษตรศาสตร์
\par เรามีทุกอย่างอยู่แล้ว — ทั้งยุทธศาสตร์ ทั้งทรัพยากร\\ แต่สิ่งที่ยังขาด คือการ \textbf{\textcolor{highlight}{“ออกแบบให้ทั้งระบบทำงานร่วมกัน”}}
\par ผมจึงเสนอกรอบคิดว่า\\ มหาวิทยาลัยเกษตรศาสตร์ต้องก้าวสู่การเป็น\\ \textbf{\textcolor{highlight}{“Life Systems University”}}
\par ไม่ใช่การเปลี่ยนตัวตน\\ แต่เป็นการ \textbf{\textcolor{highlight}{“ยกระดับจุดแข็งเดิม”}} ให้ตอบโลกที่ซับซ้อนขึ้น
\par และการจะไปถึงจุดนั้นได้\\ มี 3 เรื่องที่ต้องถูกออกแบบให้เชื่อมกันตั้งแต่ต้น
\par \textbf{\textcolor{highlight}{Governance}} — ระบบตัดสินใจต้องเร็วและโปร่งใส\\ \textbf{\textcolor{highlight}{Capability}} — คนต้องเก่งในสิ่งที่โลกต้องการจริง\\ \textbf{\textcolor{highlight}{Sustainability}} — ต้องยืนได้ในระยะยาว ทั้งวิชาการและการเงิน
\par ถ้า 3 อย่างนี้ยังแยกกันอยู่\\ มหาวิทยาลัยจะ “ขยับ” แต่ไม่ “เปลี่ยน”
\par แต่ถ้าเราออกแบบให้มันทำงานเป็นระบบเดียวกัน\\ มหาวิทยาลัยจะ \textbf{\textcolor{highlight}{“เปลี่ยนระดับ”}} ได้จริง
\par นี่คือกรอบคิดหลัก\\ และเป็นจุดตั้งต้นของทุกสิ่งที่ผมจะนำเสนอต่อไปครับ
\end{multicols}
\end{tcolorbox}

\clearpage

\section*{Slide 2: Strategic Tension}
\noindent\textcolor{muted}{Slide 2 \textbar\ Slide2.md}
\vspace{0.55em}
\begin{center}
\includegraphics[width=\linewidth,height=0.340\textheight,keepaspectratio]{/mnt/c/Users/punpi/OneDrive/KU/president/vision/qna/presentation-narrator/assets/presentation\_slides/Slide2.JPG}
\end{center}
\vspace{0.4em}
\begin{tcolorbox}[breakable,colback=panel,colframe=line,boxrule=0.5pt,arc=2mm,left=3mm,right=3mm,top=2mm,bottom=2mm]
\textbf{Version To Actually Say}
\medskip
\setlength{\columnsep}{1.1em}
\raggedcolumns
\begin{multicols}{2}
\RaggedRight
\setlength{\parskip}{0.7em}
\par สไลด์นี้ผมอยากตั้ง \textbf{\textcolor{highlight}{“โจทย์ของเวลา”}} ให้ชัดก่อนครับ
\par มหาวิทยาลัยเกษตรศาสตร์ไม่ได้อยู่ในโลกเดิมอีกต่อไป\\ แต่สิ่งที่อันตรายคือ…หลายอย่างในระบบเรา ยังคิดแบบโลกเดิม
\par ด้านหนึ่ง เรามีเป้าหมายชัด\\ ทั้ง SDGs, ranking และ digital readiness
\par แต่อีกด้านหนึ่ง โลกกำลังเร่งขึ้นพร้อมกันหลายแรง
\par \textbf{\textcolor{highlight}{AI ไม่ใช่เครื่องมืออีกต่อไป — แต่มันคือโครงสร้างพื้นฐาน}}\\ \textbf{\textcolor{highlight}{โครงสร้างประชากรกำลังเปลี่ยน — นิสิตลด รายได้ผันผวน}}\\ \textbf{\textcolor{highlight}{เศรษฐกิจใหม่ — ความเชื่อถือสำคัญกว่าปริมาณ}}
\par ปัญหาคือ\\ ถ้าเรามองสิ่งเหล่านี้เป็น \textbf{\textcolor{highlight}{“แรงภายนอก”}}
\par เราจะวางแผนได้สวย\\ แต่ระบบจริงจะ “เปราะ” ขึ้นเรื่อย ๆ
\par เพราะฉะนั้นโจทย์ของเราจริง ๆ ไม่ใช่\\ \textbf{\textcolor{highlight}{“จะเร่งหรือไม่เร่ง”}}
\par แต่คือ\\ เราจะเร่งอย่างไร…โดยที่ระบบยังมั่นคง และไม่พังระหว่างทาง
\par นี่คือสิ่งที่ผมเรียกว่า\\ \textbf{\textcolor{highlight}{strategic tension}}
\par และคำตอบของ tension นี้\\ คือสิ่งที่ผมจะพาไปในสไลด์ถัดไปครับ
\end{multicols}
\end{tcolorbox}

\clearpage

\section*{Slide 3: Identity Statement}
\noindent\textcolor{muted}{Slide 3 \textbar\ Slide3.md}
\vspace{0.55em}
\begin{center}

\includegraphics[width=\linewidth,height=0.340\textheight,keepaspectratio]{/mnt/c/Users/punpi/OneDrive/KU/president/vision/qna/presentation-narrator/assets/presentation\_slides/Slide3.JPG}
\end{center}
\vspace{0.4em}
\begin{tcolorbox}[breakable,colback=panel,colframe=line,boxrule=0.5pt,arc=2mm,left=3mm,right=3mm,top=2mm,bottom=2mm]
\textbf{Version To Actually Say}
\medskip
\setlength{\columnsep}{1.1em}
\raggedcolumns
\begin{multicols}{2}
\RaggedRight
\setlength{\parskip}{0.7em}
\par สไลด์หน้านี้ ผมไม่ได้ต้องการ\\ rebrand มหาวิทยาลัยเกษตรศาสตร์ครับ
\par ผมต้องการ \textbf{\textcolor{highlight}{“ประกาศตัวตน”}} ให้ชัด
\par Kasetsart University\\ คือ \textbf{\textcolor{highlight}{The Life Systems University}}
\par นี่ไม่ใช่คำใหม่\\ แต่เป็นการ \textbf{\textcolor{highlight}{“ขยายความหมายของรากเดิม”}}\\ ให้ตรงกับโลกปัจจุบัน
\par เพราะสิ่งที่เราทำมาตลอด\\ ไม่เคยเป็นแค่ “เกษตร”
\par แต่มันคือ\\ อาหาร สิ่งแวดล้อม เทคโนโลยี การพัฒนาคน\\ และวันนี้…รวมถึงสุขภาพของสังคม
\par ถ้ามองประเทศไทยในระยะถัดไป\\ โจทย์สำคัญไม่ใช่แค่ “แก้ปัญหา”
\par แต่คือ\\ ออกแบบระบบที่ทำให้คนไทยกินดี อยู่ดี ตั้งแต่ต้น
\par และนี่คือสิ่งที่มหาวิทยาลัยเกษตรศาสตร์\\ มี \textbf{\textcolor{highlight}{“ความพร้อมแบบหายาก”}}
\par เราไม่ได้มีแค่ความรู้\\ แต่มี ภาคสนาม งานวิจัย และเครือข่าย\\ ที่เชื่อมเรื่องเหล่านี้เข้าด้วยกันได้จริง
\par ดังนั้น The Life Systems University\\ ไม่ใช่คำสวย ๆ
\par แต่มันคือการบอกว่า\\ เราจะไม่เป็นมหาวิทยาลัยที่เก่งเป็นเรื่อง ๆ อีกต่อไป
\par แต่จะเป็น\\ มหาวิทยาลัยที่ \textbf{\textcolor{highlight}{“ออกแบบระบบชีวิตของประเทศ”}} ได้ทั้งระบบ
\par ในกรอบนี้\\ เทคโนโลยี รวมถึง AI เป็นเพียง enabler
\par ไม่ใช่ตัวตนของเรา
\par และที่สำคัญที่สุด
\par ถ้าเราไม่นิยามตัวเองให้ชัด\\ คนอื่นจะนิยามเราแทน\\ — และอาจพาเราไปผิดทาง
\par นี่คือ \textbf{\textcolor{highlight}{identity}}\\ ของมหาวิทยาลัยเกษตรศาสตร์ในระยะถัดไปครับ
\end{multicols}
\end{tcolorbox}

\clearpage

\section*{Slide 4: Technology Governance}
\noindent\textcolor{muted}{Slide 4 \textbar\ Slide4.md}
\vspace{0.55em}
\begin{center}
\includegraphics[width=\linewidth,height=0.340\textheight,keepaspectratio]{/mnt/c/Users/punpi/OneDrive/KU/president/vision/qna/presentation-narrator/assets/presentation\_slides/Slide4.JPG}
\end{center}
\vspace{0.4em}
\begin{tcolorbox}[breakable,colback=panel,colframe=line,boxrule=0.5pt,arc=2mm,left=3mm,right=3mm,top=2mm,bottom=2mm]
\textbf{Version To Actually Say}
\medskip
\setlength{\columnsep}{1.1em}
\raggedcolumns
\begin{multicols}{2}
\RaggedRight
\setlength{\parskip}{0.7em}
\par สไลด์หน้านี้ผมอยากยืนยันหลักการสำคัญข้อหนึ่งครับ
\par เทคโนโลยีของมหาวิทยาลัย\\ \textbf{\textcolor{highlight}{ต้องอยู่ใต้ธรรมาภิบาล และปัญญามนุษย์}}
\par AI และดิจิทัล ไม่ใช่เป้าหมาย\\ แต่เป็น \textbf{\textcolor{highlight}{โครงสร้างพื้นฐานของการตัดสินใจ}}
\par ถ้าฐานนี้ไม่ดี\\ สิ่งที่อยู่ข้างบนจะผิดทั้งหมด
\par ธรรมาภิบาลในที่นี้\\ ไม่ได้หมายถึงการหวังพึ่ง \textbf{\textcolor{highlight}{“คนดี”}}
\par แต่คือ\\ การออกแบบระบบให้ข้อมูลชัด โปร่งใส\\ และตรวจสอบได้ตั้งแต่ต้น
\par ไม่ใช่ค่อยวิ่งหาข้อมูล…เวลาจะมีประชุม
\par เมื่อข้อมูลอยู่ในระบบเดียวกัน\\ ผู้บริหารและสภา จะตัดสินใจจาก \textbf{\textcolor{highlight}{หลักฐานเดียวกัน}}
\par ไม่ใช่คนละชุด คนละความจริง
\par ตรงนี้ AI เข้ามามีบทบาท
\par \textbf{\textcolor{highlight}{ไม่ใช่เพื่อ “ตัดสินใจแทน”}}\\ แต่เพื่อ\\ ช่วยแยกแยะ จัดกลุ่ม และทำให้เราเห็น pattern ที่ซ่อนอยู่ในข้อมูล
\par แต่สุดท้าย\\ การตัดสินใจยังเป็นของมนุษย์\\ และต้องมีความรับผิดชอบเชิงนโยบาย
\par ถ้าระบบนี้เกิดขึ้นจริง
\par สิ่งที่เปลี่ยนจะไม่ใช่แค่ KPI ดีขึ้น
\par แต่คือ\\ ความขัดแย้งในองค์กรจะลดลง\\ เพราะทุกคนทำงานบนข้อมูลชุดเดียวกัน
\par และทั้งมหาวิทยาลัย\\ จะค่อย ๆ เชื่อมเข้าสู่ \textbf{\textcolor{highlight}{data stream เดียว}}
\par ตั้งแต่งานวิจัย การเรียนการสอน\\ ไปจนถึงการบริหารและการตัดสินใจระดับสภา
\par นี่คือ \textbf{\textcolor{highlight}{“โครงสร้างฐาน”}}\\ ที่ทำให้ Life Systems University เกิดขึ้นได้จริงครับ
\end{multicols}
\end{tcolorbox}

\clearpage

\section*{Slide 5: Life Systems}
\noindent\textcolor{muted}{Slide 5 \textbar\ Slide5.md}
\vspace{0.55em}
\begin{center}
\includegraphics[width=\linewidth,height=0.340\textheight,keepaspectratio]{/mnt/c/Users/punpi/OneDrive/KU/president/vision/qna/presentation-narrator/assets/presentation\_slides/Slide5.JPG}
\end{center}
\vspace{0.4em}
\begin{tcolorbox}[breakable,colback=panel,colframe=line,boxrule=0.5pt,arc=2mm,left=3mm,right=3mm,top=2mm,bottom=2mm]
\textbf{Version To Actually Say}
\medskip
\setlength{\columnsep}{1.1em}
\raggedcolumns
\begin{multicols}{2}
\RaggedRight
\setlength{\parskip}{0.7em}
\par สไลด์หน้านี้ ผมอยากให้เห็นภาพง่าย ๆ ครับ
\par มหาวิทยาลัยเกษตรศาสตร์\\ มี \textbf{\textcolor{highlight}{“ระบบชีวิต”}} นี้อยู่แล้ว
\par ตั้งแต่\\ \textbf{\textcolor{highlight}{agriculture → food → health → sustainability}}
\par ปัญหาคือที่ผ่านมา\\ สิ่งเหล่านี้ \textbf{\textcolor{highlight}{“อยู่ครบ”}}\\ แต่ยังไม่ได้ \textbf{\textcolor{highlight}{“เชื่อมเป็นระบบเดียว”}}
\par \textbf{\textcolor{highlight}{Life Systems University}}\\ จึงไม่ใช่การตั้งชื่อใหม่
\par แต่คือ\\ การจัดวางสิ่งที่เรามี ให้ทำงานเป็นระบบเดียวกัน\\ อย่างมีทิศทาง
\par เมื่อมองแบบนี้\\ โจทย์ของประเทศจะชัดขึ้นทันที
\par ไม่ใช่แค่เรื่องใดเรื่องหนึ่ง\\ แต่คือโจทย์ระดับระบบ เช่น
\par climate adaptation\\ food security\\ one health\\ bio innovation
\par และนี่คือจุดที่ผมคิดว่า\\ KU มีสิทธิ์เป็นผู้นำ
\par ไม่ใช่เพราะเราทำเรื่องใดเรื่องหนึ่งเก่งที่สุด
\par แต่เพราะ\\ เรามี \textbf{\textcolor{highlight}{“ความต่อเนื่องของระบบ”}} ตั้งแต่ต้นน้ำถึงปลายน้ำ
\par ซึ่งเป็นสิ่งที่มหาวิทยาลัยส่วนใหญ่ “ไม่มี”
\par แน่นอน\\ การเชื่อมระบบนี้ให้ทำงานจริง\\ ต้องมี technology backbone ที่แข็งแรง
\par ทั้ง AI, data, sensors, genomics, digital twin
\par แต่ผมอยากเน้นว่า
\par เทคโนโลยีไม่ใช่สารหลักของสไลด์นี้
\par สารหลักคือ\\ KU ต้องเป็นที่ที่ความรู้เรื่อง \textbf{\textcolor{highlight}{“ระบบชีวิต” ถูกเชื่อมเข้าด้วยกัน}}\\ และแปลงเป็น impact ต่อประเทศได้จริง
\end{multicols}
\end{tcolorbox}

\clearpage

\section*{Slide 6: National Role}
\noindent\textcolor{muted}{Slide 6 \textbar\ Slide6.md}
\vspace{0.55em}
\begin{center}
\includegraphics[width=\linewidth,height=0.340\textheight,keepaspectratio]{/mnt/c/Users/punpi/OneDrive/KU/president/vision/qna/presentation-narrator/assets/presentation\_slides/Slide6.JPG}
\end{center}
\vspace{0.4em}
\begin{tcolorbox}[breakable,colback=panel,colframe=line,boxrule=0.5pt,arc=2mm,left=3mm,right=3mm,top=2mm,bottom=2mm]
\textbf{Version To Actually Say}
\medskip
\setlength{\columnsep}{1.1em}
\raggedcolumns
\begin{multicols}{2}
\RaggedRight
\setlength{\parskip}{0.7em}
\par สไลด์หน้านี้\\ ผมอยากชวนมองบทบาทของมหาวิทยาลัย\\ ในระดับประเทศให้ชัดครับ
\par มหาวิทยาลัยไม่จำเป็นต้องทำทุกอย่างเอง\\ แต่ต้องทำในสิ่งที่ \textbf{\textcolor{highlight}{“ระบบอื่นทำไม่ได้”}}
\par โดยเฉพาะ 3 เรื่อง
\par \textbf{\textcolor{highlight}{คนที่ยังไม่พร้อม}}\\ \textbf{\textcolor{highlight}{สิ่งที่ประเทศยังขาด}}\\ และ\textbf{\textcolor{highlight}{งานที่ตลาดไม่ทำ}}
\par นี่คือเหตุผลที่ผมใช้คำว่า
\par \textbf{\textcolor{highlight}{Trusted System Partner for Capability Building}}
\par บทบาทของมหาวิทยาลัย\\ คือช่วย \textbf{\textcolor{highlight}{“ลดความเสี่ยงเชิงระบบของประเทศ”}}
\par ทำงานพื้นฐานที่ไม่มีใครทำ\\ รักษาฐานความรู้ระยะยาว\\ และเปิดพื้นที่ให้รัฐ อุตสาหกรรม และสังคม\\ ทำงานร่วมกันได้จริง
\par เพราะถ้าปล่อยให้ทุกอย่างขับเคลื่อน\\ ด้วย logic ของตลาดอย่างเดียว
\par งานที่สร้าง capability ระยะยาว\\ จำนวนมาก…จะไม่เกิด
\par ในความหมายนี้\\ มหาวิทยาลัยไม่ควรลงไปเป็น operator ในทุกเรื่อง
\par เพราะไม่ใช่ความถนัด\\ ไม่ใช่พันธกิจ\\ และอาจทำให้ระบบเสียสมดุล
\par บทบาทที่เหมาะกว่าคือ
\par \textbf{\textcolor{highlight}{ออกแบบโจทย์}}\\ \textbf{\textcolor{highlight}{พัฒนาคน}}\\ \textbf{\textcolor{highlight}{สร้างความรู้}}\\ และเชื่อมทุกฝ่ายให้\textbf{\textcolor{highlight}{ทำงานเป็นระบบเดียวกัน}}
\par ความร่วมมือกับภายนอกก็เช่นกัน
\par เราต้องเป็น trusted partner จริง\\ ไม่ใช่รับทุกดีล\\ ไม่ใช่ทำ MOU เชิงพิธีการ
\par แต่เลือกเฉพาะความร่วมมือ\\ ที่ \textbf{\textcolor{highlight}{“สร้างคน”}} และ \textbf{\textcolor{highlight}{“สร้างความรู้”}} ได้จริง
\par สุดท้ายแล้ว
\par ความน่าเชื่อถือของมหาวิทยาลัย\\ ไม่ได้มาจากการทำทุกอย่าง
\par แต่มาจาก\\ การรู้ว่าตัวเองต้องทำอะไร — และทำสิ่งนั้นให้ประเทศเชื่อถือได้
\end{multicols}
\end{tcolorbox}

\clearpage

\section*{Slide 7: Outcome-Driven Education}
\noindent\textcolor{muted}{Slide 7 \textbar\ Slide7.md}
\vspace{0.55em}
\begin{center}
\includegraphics[width=\linewidth,height=0.340\textheight,keepaspectratio]{/mnt/c/Users/punpi/OneDrive/KU/president/vision/qna/presentation-narrator/assets/presentation\_slides/Slide7.JPG}
\end{center}
\vspace{0.4em}
\begin{tcolorbox}[breakable,colback=panel,colframe=line,boxrule=0.5pt,arc=2mm,left=3mm,right=3mm,top=2mm,bottom=2mm]
\textbf{Version To Actually Say}
\medskip
\setlength{\columnsep}{1.1em}
\raggedcolumns
\begin{multicols}{2}
\RaggedRight
\setlength{\parskip}{0.7em}
\par สไลด์หน้านี้ ผมอยากทำให้คำว่า\\ \textbf{\textcolor{highlight}{“บัณฑิตที่ทำงานได้จริง”}} เป็นระบบที่เกิดขึ้นได้จริงครับ
\par ถ้าเราต้องการบัณฑิตที่\\ ทำงานได้ เรียนรู้ต่อเนื่อง และสร้างคุณค่าให้สังคม
\par เราไม่สามารถปล่อยให้ภาควิชารับภาระนี้อยู่ลำพังได้
\par ภาควิชา\\ ยังต้องเป็น \textbf{\textcolor{highlight}{“เจ้าของการเรียนรู้”}} อย่างเต็มที่\\ และต้องมีอิสระทางวิชาการ
\par แต่บทบาทของมหาวิทยาลัยต้องเปลี่ยน
\par จาก “ผู้กำกับหลักสูตร”\\ มาเป็น\\ \textbf{\textcolor{highlight}{“ระบบสนับสนุนที่ทำให้หลักสูตรเชื่อมกับโลกจริงได้”}}
\par โดยไม่เพิ่มขั้นตอน\\ และไม่ดึงอำนาจจากภาควิชา
\par กลไกที่ผมจะทำ มี 4 ส่วน
\par หนึ่ง — \textbf{\textcolor{highlight}{internship เชิงเป้าหมาย}}\\ ไม่ใช่แค่ไปฝึกงาน แต่เชื่อมกับโจทย์ที่หลักสูตรต้องการ
\par สอง — \textbf{\textcolor{highlight}{career tracking}}\\ เราต้องรู้ว่าบัณฑิตออกไปแล้วเป็นอย่างไรจริง ไม่ใช่เดา
\par สาม — \textbf{\textcolor{highlight}{data feedback}}\\ ข้อมูลจากโลกจริง ต้องย้อนกลับมาปรับหลักสูตรได้
\par สี่ — \textbf{\textcolor{highlight}{funding support}}\\ ให้ภาควิชามีทรัพยากรทดลองสิ่งใหม่ โดยไม่ผลักภาระให้อาจารย์
\par ในระบบนี้\\ บทบาทต้องชัด
\par ภาควิชา — ออกแบบการเรียนรู้\\ มหาวิทยาลัย — สร้างระบบสนับสนุน\\ สังคม/อุตสาหกรรม — สะท้อนสมรรถนะที่จำเป็น
\par ไม่ใช่เข้ามากำหนดเนื้อหา
\par ถ้าทำแบบนี้ได้
\par \textbf{\textcolor{highlight}{Outcome-driven education}}\\ จะไม่ใช่การสอนตามตำแหน่งงาน
\par แต่จะเป็น\\ การยกระดับคุณภาพหลักสูตรด้วยข้อมูลจริงและโจทย์จริง
\par และนี่คือวิธีที่มหาวิทยาลัยทำหน้าที่\\ \textbf{\textcolor{highlight}{Trusted System Partner for Capability Building}}\\ ผ่าน \textbf{\textcolor{highlight}{“การพัฒนาคน”}} อย่างเป็นระบบ
\end{multicols}
\end{tcolorbox}

\clearpage

\section*{Slide 8: Human Sustainability}
\noindent\textcolor{muted}{Slide 8 \textbar\ Slide8.md}
\vspace{0.55em}
\begin{center}
\includegraphics[width=\linewidth,height=0.340\textheight,keepaspectratio]{/mnt/c/Users/punpi/OneDrive/KU/president/vision/qna/presentation-narrator/assets/presentation\_slides/Slide8.JPG}
\end{center}
\vspace{0.4em}
\begin{tcolorbox}[breakable,colback=panel,colframe=line,boxrule=0.5pt,arc=2mm,left=3mm,right=3mm,top=2mm,bottom=2mm]
\textbf{Version To Actually Say}
\medskip
\setlength{\columnsep}{1.1em}
\raggedcolumns
\begin{multicols}{2}
\RaggedRight
\setlength{\parskip}{0.7em}
\par สไลด์หน้านี้ ผมอยากย้ำเรื่องที่สำคัญมากครับ
\par ถ้าคนในมหาวิทยาลัยอยู่ไม่ไหว\\ มหาวิทยาลัยจะไม่มีวันยั่งยืน
\par เพราะฉะนั้น \textbf{\textcolor{highlight}{Human Sustainability}}\\ ไม่ใช่เรื่อง soft\\ และไม่ใช่สวัสดิการปลายทาง
\par แต่มันคือ\\ \textbf{\textcolor{highlight}{โครงสร้างพื้นฐานของสถาบัน}}
\par ปัญหาที่เราเห็นวันนี้\\ ไม่ใช่แค่คนเหนื่อย
\par แต่คือ\\ หลายระบบยังพึ่ง \textbf{\textcolor{highlight}{“ความเสียสละ”}} มากกว่า \textbf{\textcolor{highlight}{“การออกแบบ”}}
\par อาจารย์ต้องสอน ดูแลนิสิต ทำวิจัย\\ และขับเคลื่อนองค์กร
\par แต่ระบบความก้าวหน้าและแรงจูงใจ\\ ไม่ได้สะท้อนภารกิจเหล่านี้อย่างสมดุล
\par สิ่งที่ควรเป็น \textbf{\textcolor{highlight}{“งานหลัก”}}\\ กลับต้องอาศัยการเสียสละ
\par ในขณะที่สิ่งที่ \textbf{\textcolor{highlight}{“วัดง่าย”}}\\ กลับได้รางวัลชัดกว่า
\par นี่คือเหตุผลที่คนค่อย ๆ หมดแรง
\par เพราะฉะนั้น สิ่งที่ต้องเปลี่ยน\\ ไม่ใช่ให้คนอดทนมากขึ้น
\par แต่คือ\\ \textbf{\textcolor{highlight}{เปลี่ยนระบบให้สะท้อนงานจริงของมหาวิทยาลัย}}
\par ทั้งในเรื่อง\\ การประเมิน ความก้าวหน้า และแรงจูงใจ
\par ควบคู่กับการทำเรื่องพื้นฐานให้ชัด
\par ลดงานที่ไม่จำเป็น\\ มีระบบรับฟังที่ไม่เงียบหาย\\ มีความชัดเจนเรื่องสิทธิและสวัสดิการ\\ และสภาพแวดล้อมที่ทำงานได้จริง
\par และที่สำคัญ
\par ระบบจากศูนย์กลาง\\ ต้องตั้งอยู่บนความเข้าใจว่า\\ หน่วยย่อยไม่เหมือนกัน
\par เพราะคณะและภาควิชา\\ คือหัวใจของมหาวิทยาลัย
\par สุดท้ายแล้ว
\par Human Sustainability\\ ไม่ใช่การดูแลคนให้ \textbf{\textcolor{highlight}{“อยู่ได้”}}
\par แต่คือ\\ การทำให้คนทำงานหลักของมหาวิทยาลัยได้อย่างยั่งยืน\\ โดยไม่ต้องใช้ความเสียสละเป็นเชื้อเพลิงหลัก
\end{multicols}
\end{tcolorbox}

\clearpage

\section*{Slide 9: Trust by Design}
\noindent\textcolor{muted}{Slide 9 \textbar\ Slide9.md}
\vspace{0.55em}
\begin{center}
\includegraphics[width=\linewidth,height=0.340\textheight,keepaspectratio]{/mnt/c/Users/punpi/OneDrive/KU/president/vision/qna/presentation-narrator/assets/presentation\_slides/Slide9.JPG}
\end{center}
\vspace{0.4em}
\begin{tcolorbox}[breakable,colback=panel,colframe=line,boxrule=0.5pt,arc=2mm,left=3mm,right=3mm,top=2mm,bottom=2mm]
\textbf{Version To Actually Say}
\medskip
\setlength{\columnsep}{1.1em}
\raggedcolumns
\begin{multicols}{2}
\RaggedRight
\setlength{\parskip}{0.7em}
\par สไลด์หน้านี้ ผมพูดตรงไปตรงมาครับ
\par ความเชื่อมั่นของสถาบัน\\ ต้องไม่ขึ้นกับว่าเราโชคดีได้คนแบบไหน
\par ถ้ายังต้องหวังว่าได้ \textbf{\textcolor{highlight}{“คนดี”}} แล้วทุกอย่างจะดี\\ แปลว่าเรายังออกแบบระบบไม่ดีพอ
\par โจทย์ของ governance จึงไม่ใช่แค่ควบคุม\\ แต่คือ\\ ออกแบบให้การตัดสินใจ \textbf{\textcolor{highlight}{“สม่ำเสมอ โปร่งใส และตรวจสอบได้”}}\\ แม้ตัวคนจะเปลี่ยน
\par ผมไม่ได้บอกว่าไม่ต้องมีดุลยพินิจ
\par แต่ดุลยพินิจต้องอยู่ใน \textbf{\textcolor{highlight}{กรอบที่ชัด}}\\ ไม่ใช่ปล่อยให้ discretion กลายเป็นตัวแทนของระบบ
\par ปัญหาที่เห็นชัดในมหาวิทยาลัยคือ\\ บทบาทไม่ชัด อำนาจไม่ชัด และความรับผิดชอบไม่ชัด
\par เราใช้คำว่า “กรรมการ” กว้างเกินไป\\ ทั้งที่จริงมีทั้งคนที่ “ตัดสิน”\\ คนที่ “แทน”\\ และคนที่ “ให้คำปรึกษา”
\par ถ้าบทบาทไม่ชัด\\ สุดท้าย\\ ความรับผิดชอบจะพร่า และดุลยพินิจจะไหลเข้าแทนระบบ
\par สิ่งที่ผมจะทำ มี 4 เรื่องชัด
\par หนึ่ง — \textbf{\textcolor{highlight}{Standardized Governance}}\\ กำหนดบทบาท สายการตัดสินใจ และ workflow ให้ชัด\\ ลดความผันผวนตามตัวบุคคล
\par สอง — \textbf{\textcolor{highlight}{Reduced Discretion (where it shouldn’t be)}}\\ ลดดุลยพินิจในเรื่องที่ควรเป็นระบบ\\ เพื่อปิดช่องความเสี่ยงเชิงโครงสร้าง
\par สาม — \textbf{\textcolor{highlight}{Transparent \& Secure Evaluation}}\\ แยกระบบประเมินออกจากผู้ถูกประเมิน\\ ข้อมูลปลอดภัย และตรวจสอบได้จริง
\par สี่ — \textbf{\textcolor{highlight}{Council Data Dashboard}}\\ ให้สภากำกับจาก \textbf{\textcolor{highlight}{“ข้อมูลชุดเดียวกัน”}}\\ เห็นภาพรวมและความเสี่ยง\\ โดยไม่ต้องลงไปแทรกแซงรายกรณี
\par ถ้าระบบนี้เกิดขึ้นจริง
\par ความเป็นธรรม ความต่อเนื่อง และความไว้วางใจ\\ จะไม่ต้องผูกกับตัวบุคคลอีกต่อไป
\par แต่จะถูก \textbf{\textcolor{highlight}{“ออกแบบไว้ในระบบ”}} ตั้งแต่ต้น
\end{multicols}
\end{tcolorbox}

\clearpage

\section*{Slide 10: Academic Engine}
\noindent\textcolor{muted}{Slide 10 \textbar\ Slide10.md}
\vspace{0.55em}
\begin{center}
\includegraphics[width=\linewidth,height=0.210\textheight,keepaspectratio]{/mnt/c/Users/punpi/OneDrive/KU/president/vision/qna/presentation-narrator/assets/presentation\_slides/Slide10.JPG}
\end{center}
\vspace{0.4em}
\begin{tcolorbox}[breakable,colback=panel,colframe=line,boxrule=0.5pt,arc=2mm,left=3mm,right=3mm,top=2mm,bottom=2mm]
\textbf{Version To Actually Say}
\medskip
\setlength{\columnsep}{1.1em}
\raggedcolumns
\begin{multicols}{2}
\RaggedRight
\setlength{\parskip}{0.7em}
\par สไลด์หน้านี้ ผมอยากชัดเจนเรื่องหนึ่งครับ
\par ถ้า KU จะสร้าง impact ได้มากขึ้น\\ เราต้องไม่เลือกระหว่าง \textbf{\textcolor{highlight}{“ฐาน”}} กับ \textbf{\textcolor{highlight}{“อนาคต”}}
\par แต่ต้องทำให้มันทำงานร่วมกัน
\par \textbf{\textcolor{highlight}{Strong foundations × }}dynamic frontiers====
\par ถ้าฐานไม่แข็ง\\ frontier จะเป็นแค่คำใหม่
\par แต่ถ้ามีแต่ฐาน\\ โดยไม่เชื่อมกับโจทย์โลก
\par ศักยภาพนั้นจะไม่ถูกแปลงเป็น impact
\par คำว่า dynamic frontiers\\ ไม่ได้หมายถึงการตั้งหน่วยงานใหม่ไปเรื่อย ๆ
\par แต่หมายถึง\\ การเอาจุดแข็งเดิมของ KU ไปเชื่อมกับโจทย์ที่กำลังเร่งตัว
\par เช่น\\ AI \& data\\ sustainability\\ food \& bio systems\\ digital economy\\ policy \& governance
\par สิ่งสำคัญคือ
\par frontier ไม่ใช่ของใหม่ที่ลอยอยู่ข้างบน\\ แต่เป็น \textbf{\textcolor{highlight}{“ชั้น”}} ที่ทำให้ core เดิม\\ เข้าไปสร้างคุณค่าในโลกใหม่ได้
\par และในอีกมุมหนึ่ง
\par core ของแต่ละสาขา\\ ต้องสามารถเป็น \textbf{\textcolor{highlight}{“patch”}} ให้กันและกันได้
\par นี่คือจุดที่ยากที่สุดในมหาวิทยาลัย
\par เพราะวันนี้\\ เราทำ MOU กับคนนอกได้ง่าย
\par แต่ทำความร่วมมือ \textbf{\textcolor{highlight}{“ข้ามภาควิชา”}} ให้เกิดจริง\\ กลับยากกว่า
\par เพราะฉะนั้น สิ่งที่ผมจะทำ\\ ไม่ใช่แค่บอกให้ร่วมมือกัน
\par แต่ต้องมี \textbf{\textcolor{highlight}{mechanism}} ที่ทำให้มันเกิด
\par ทั้งในรูปของ
\par วิชาและหลักสูตรข้ามศาสตร์\\ โจทย์วิจัยร่วม\\ และ funding ที่สนับสนุนการทำงานข้ามสาขา
\par ตัวอย่างที่เริ่มเห็นแล้วคือ \textbf{\textcolor{highlight}{AIEP}}
\par ซึ่งไม่ได้เอา AI ไปแปะ\\ แต่เอา AI ไป \textbf{\textcolor{highlight}{“ยืนบน domain knowledge จริง”}}
\par เด็กยังต้องแข็งในสาขาของตัวเอง\\ แล้วค่อยต่อยอด
\par และที่สำคัญ\\ มันพิสูจน์ว่า\\ ถ้ามีเจ้าภาพ และมีระบบรองรับ\\ core ของสาขาหนึ่ง สามารถเป็น patch ให้สาขาอื่นได้จริง
\par สุดท้ายแล้ว
\par KU ไม่ได้ต้องเลิกเป็นมหาวิทยาลัยตามสาขา
\par แต่ต้อง\\ ใช้ความแข็งแรงในสาขา\\ แล้วจัดพลังใหม่ตามโจทย์โลก
\par เพื่อให้เกิด\\ วิชาใหม่ หลักสูตรใหม่\\ และงานวิจัยที่สร้าง impact ได้จริง
\end{multicols}
\end{tcolorbox}

\clearpage

\section*{Slide 11: Digital Governance}
\noindent\textcolor{muted}{Slide 11 \textbar\ Slide11.md}
\vspace{0.55em}
\begin{center}
\includegraphics[width=\linewidth,height=0.210\textheight,keepaspectratio]{/mnt/c/Users/punpi/OneDrive/KU/president/vision/qna/presentation-narrator/assets/presentation\_slides/Slide11.JPG}
\end{center}
\vspace{0.4em}
\begin{tcolorbox}[breakable,colback=panel,colframe=line,boxrule=0.5pt,arc=2mm,left=3mm,right=3mm,top=2mm,bottom=2mm]
\textbf{Version To Actually Say}
\medskip
\setlength{\columnsep}{1.1em}
\raggedcolumns
\begin{multicols}{2}
\RaggedRight
\setlength{\parskip}{0.7em}
\par สไลด์หน้านี้ ผมอยากชัดเจนเรื่องหนึ่งครับ
\par ธรรมาภิบาลที่ดี\\ ต้องตั้งอยู่บน \textbf{\textcolor{highlight}{“ข้อมูลชุดเดียวกัน”}} ที่ทุกคนเห็นร่วมกัน
\par ถ้าฝ่ายบริหารเห็นตัวเลขชุดหนึ่ง\\ สภาเห็นอีกชุดหนึ่ง\\ และส่วนงานเห็นอีกชุดหนึ่ง
\par ต่อให้เราพูดเรื่องความโปร่งใสแค่ไหน
\par การตัดสินใจจะกลับไปอยู่ที่ความรู้สึก\\ อำนาจ\\ หรือการตีความเฉพาะหน้า
\par เพราะฉะนั้น หน้านี้ไม่ใช่ IT project
\par แต่คือ\\ \textbf{\textcolor{highlight}{decision infrastructure}} ของมหาวิทยาลัย
\par ผมมองเป็น 3 ชั้น
\par \textbf{\textcolor{highlight}{DATA}}\\ เราต้องมีโครงสร้างข้อมูลกลาง\\ ที่เชื่อมกันทั้งระบบ\\ รู้ว่าใครเป็นเจ้าของข้อมูล\\ และข้อมูลไหนใช้ตอบคำถามอะไร
\par \textbf{\textcolor{highlight}{DECISION}}\\ ข้อมูลต้องไม่หยุดที่รายงาน\\ แต่ต้องเข้าไปอยู่ใน\\ การจัดสรรทรัพยากร และการตัดสินใจจริง
\par \textbf{\textcolor{highlight}{INTELLIGENCE}}\\ ข้อมูลต้องช่วยให้เกิด \textbf{\textcolor{highlight}{“ความเป็นธรรม”}}
\par ไม่ใช่แค่ดูตัวเลข\\ แต่ต้องเห็น\\ ภาระงาน บริบท และข้อจำกัดของแต่ละส่วนงาน
\par ก่อนจะตัดสิน
\par อีกเรื่องที่สำคัญมากคือ
\par ความโปร่งใสที่ดี\\ ต้องไม่ทำให้คนกลัวข้อมูล
\par ถ้าข้อมูลถูกใช้เพื่อ \textbf{\textcolor{highlight}{“จับผิด”}}\\ ระบบจะไม่เกิด
\par แต่ถ้าข้อมูลถูกใช้เพื่อ \textbf{\textcolor{highlight}{“สนับสนุน”}}
\par ความไว้วางใจจะเกิดเอง
\par เช่น\\ ถ้าเห็นว่าหน่วยงานใดมีภาระงานสูง\\ สิ่งที่ตามมาควรเป็นการช่วย\\ ไม่ใช่การตำหนิ
\par และถ้าระบบทำให้\\ กรอกครั้งเดียว ใช้ได้ทั้งระบบ
\par คนจะเริ่มเชื่อระบบเร็วขึ้น
\par ปัญหาของหลายองค์กร\\ ไม่ใช่ไม่มีข้อมูล
\par แต่คือ\\ คนทำข้อมูล คนใช้ข้อมูล และคนตัดสินใจ\\ ไม่ได้อยู่ในระบบเดียวกัน
\par เพราะฉะนั้น \textbf{\textcolor{highlight}{data governance}} ที่ดี\\ ต้องทำให้ทุกฝ่าย\\ เชื่อมกันอยู่ใน \textbf{\textcolor{highlight}{“ภาษากลางเดียวกัน”}}
\par ถ้าทำได้
\par การตัดสินใจของมหาวิทยาลัย\\ จะโปร่งใส รับผิดชอบได้\\ และเป็นธรรมมากขึ้น
\par และสำหรับผม\\ data governance คือ priority ของการเปลี่ยนแปลงนี้ครับ
\end{multicols}
\end{tcolorbox}

\clearpage

\section*{Slide 12: Financial Sustainability}
\noindent\textcolor{muted}{Slide 12 \textbar\ Slide12.md}
\vspace{0.55em}
\begin{center}
\includegraphics[width=\linewidth,height=0.170\textheight,keepaspectratio]{/mnt/c/Users/punpi/OneDrive/KU/president/vision/qna/presentation-narrator/assets/presentation\_slides/Slide12.JPG}
\end{center}
\vspace{0.4em}
\begin{tcolorbox}[breakable,colback=panel,colframe=line,boxrule=0.5pt,arc=2mm,left=3mm,right=3mm,top=2mm,bottom=2mm]
\textbf{Version To Actually Say}
\medskip
\setlength{\columnsep}{1.1em}
\raggedcolumns
\begin{multicols}{2}
\RaggedRight
\setlength{\parskip}{0.7em}
\par สไลด์หน้านี้ ผมไม่ได้พูดเรื่อง “หาเงิน” อย่างเดียวครับ
\par ผมพูดเรื่อง\\ ทำอย่างไรให้มหาวิทยาลัยมี\textbf{\textcolor{highlight}{ความมั่นคงทางการเงิน}}\\ โดยไม่เสียความเป็นอิสระทางวิชาการ
\par ถ้าฐานการเงินเปราะบาง
\par มหาวิทยาลัยจะถูกบีบให้\\ ตัดสินใจระยะสั้น\\ วิ่งตามแรงกดดัน\\ และแก้ปัญหาเฉพาะหน้าไปเรื่อย ๆ
\par เพราะฉะนั้นแกนแรกคือ\\ \textbf{\textcolor{highlight}{financial discipline}}
\par ไม่ใช่แค่รัดเข็มขัด
\par แต่คือ\\ เห็นความเสี่ยงก่อนมันเกิด
\par ผมใช้หลัก\\ \textbf{\textcolor{highlight}{monitor before disaster}}
\par เราต้องเห็นแนวโน้มรายได้\\ ภาระผูกพัน\\ และจุดเสี่ยงล่วงหน้า
\par ไม่ใช่รอให้ขาดดุลแล้วค่อยแก้
\par แกนที่สองคือ\\ \textbf{\textcolor{highlight}{data-driven allocation}}
\par งบประมาณต้องสะท้อน \textbf{\textcolor{highlight}{“ยุทธศาสตร์”}}\\ ไม่ใช่ \textbf{\textcolor{highlight}{“ประวัติศาสตร์”}}
\par และการจัดสรรต้องอยู่บนข้อมูลจริง\\ ไม่ใช่ inertia หรืออำนาจต่อรอง
\par แกนที่สามคือ\\ \textbf{\textcolor{highlight}{clear business model}}
\par ไม่ได้แปลว่าต้องกำไรสูงสุด
\par แต่ต้องรู้ว่า
\par รายได้หลักเลี้ยงภารกิจหลักได้หรือไม่\\ ต้นทุนจริงอยู่ตรงไหน\\ และกิจกรรมใด \textbf{\textcolor{highlight}{“สอดคล้องกับพันธกิจ”}} จริง
\par ตรงนี้เชื่อมกับบทบาท\\ \textbf{\textcolor{highlight}{Trusted System Partner}}
\par ถ้าเราทำหน้าที่นี้ได้จริง
\par รายได้ใหม่จะเกิดจาก\\ \textbf{\textcolor{highlight}{mission-aligned growth}}
\par ไม่ใช่การวิ่งไล่รายได้ที่ไม่ใช่ตัวตน
\par ผมไม่คิดว่าคำตอบคือ\\ ขึ้นค่าเทอม\\ หรือหาเงินทุกทาง
\par และผมก็ไม่คิดว่า\\ เงินบริจาคควรเป็นเสาหลัก
\par แต่เราสามารถสร้างรายได้เสริมที่มีคุณภาพได้
\par เช่น
\par การต่อยอดทรัพย์สินทางปัญญา\\ งานวิจัยสัญญาที่มีคุณภาพ\\ บริการเฉพาะทาง\\ และกองทุนเชิงยุทธศาสตร์
\par ถ้าจะระดมทุน\\ ต้องเป็น \textbf{\textcolor{highlight}{thematic fundraising}}\\ ที่ผู้ให้รู้ว่าเงินสร้าง impact อะไร
\par สุดท้ายแล้ว
\par ความมั่นคงทางการเงิน\\ ไม่ใช่การอยู่รอดแบบประคองตัว
\par แต่คือ\\ การมี\textbf{\textcolor{highlight}{อิสระในการตัดสินใจ}}\\ และขยายภารกิจได้อย่างรับผิดชอบ\\ บนฐานข้อมูลจริง
\end{multicols}
\end{tcolorbox}

\clearpage

\section*{Slide 13: Leadership Positioning}
\noindent\textcolor{muted}{Slide 13 \textbar\ Slide13.md}
\vspace{0.55em}
\begin{center}
\includegraphics[width=\linewidth,height=0.340\textheight,keepaspectratio]{/mnt/c/Users/punpi/OneDrive/KU/president/vision/qna/presentation-narrator/assets/presentation\_slides/Slide13.JPG}
\end{center}
\vspace{0.4em}
\begin{tcolorbox}[breakable,colback=panel,colframe=line,boxrule=0.5pt,arc=2mm,left=3mm,right=3mm,top=2mm,bottom=2mm]
\textbf{Version To Actually Say}
\medskip
\setlength{\columnsep}{1.1em}
\raggedcolumns
\begin{multicols}{2}
\RaggedRight
\setlength{\parskip}{0.7em}
\par สไลด์หน้านี้ ไม่ใช่การสรุปประวัติของผมครับ
\par แต่เป็นคำตอบว่า\\ ทำไมผมจึงเหมาะกับบทบาทอธิการบดีในช่วงนี้
\par โจทย์ของมหาวิทยาลัยวันนี้\\ ไม่ใช่เรื่องใดเรื่องหนึ่ง
\par แต่เป็นการขยับ \textbf{\textcolor{highlight}{“ทั้งระบบ”}} ไปพร้อมกัน
\par เพราะฉะนั้น ผู้นำต้องทำงานได้ครบ 4 มิติ
\par หนึ่ง — \textbf{\textcolor{highlight}{มองระบบ}}\\ เข้าใจโครงสร้าง อำนาจ และกลไกจริงของมหาวิทยาลัย\\ ไม่ใช่เห็นแค่ปัญหาเฉพาะจุด
\par สอง — \textbf{\textcolor{highlight}{พัฒนาคน}}\\ มหาวิทยาลัยไม่ใช่องค์กรอย่างเดียว\\ แต่คือระบบสร้างศักยภาพคนของประเทศ
\par สาม — \textbf{\textcolor{highlight}{เชื่อมโลก}}\\ ความร่วมมือระหว่างประเทศต้องสร้าง \textbf{\textcolor{highlight}{“ผลจริง”}}\\ ไม่ใช่แค่พิธีการ
\par สี่ — \textbf{\textcolor{highlight}{สร้างของที่อยู่ต่อได้}}\\ ไม่ใช่แค่เริ่มโครงการ\\ แต่ทำให้มันกลายเป็น \textbf{\textcolor{highlight}{“โครงสร้างของมหาวิทยาลัย”}}
\par สิ่งที่ผมทำมาตลอด\\ คือทำงานใน 4 มิตินี้ \textbf{\textcolor{highlight}{“พร้อมกัน”}}
\par ไม่ว่าจะเป็น\\ งานระบบระดับมหาวิทยาลัย\\ การพัฒนาคนผ่านโอลิมปิกและหลักสูตร\\ ความร่วมมือระดับนานาชาติ\\ หรือการสร้างโครงการที่อยู่ต่อได้จริง เช่น AIEP
\par เพราะฉะนั้น จุดแข็งของผม\\ ไม่ใช่ว่าทำอะไรมาเยอะ
\par แต่คือ\\ เข้าใจว่าระบบมหาวิทยาลัยทำงานอย่างไร\\ และทำให้มัน\textbf{\textcolor{highlight}{ขยับได้จริง}}
\par และผมเชื่อว่า\\ นี่คือสิ่งที่มหาวิทยาลัยต้องการ\\ ในช่วงเปลี่ยนผ่านนี้ครับ
\end{multicols}
\end{tcolorbox}

\clearpage

\section*{Slide 14: 4-Year Execution Map}
\noindent\textcolor{muted}{Slide 14 \textbar\ Slide14.md}
\vspace{0.55em}
\begin{center}
\includegraphics[width=\linewidth,height=0.340\textheight,keepaspectratio]{/mnt/c/Users/punpi/OneDrive/KU/president/vision/qna/presentation-narrator/assets/presentation\_slides/Slide14.JPG}
\end{center}
\vspace{0.4em}
\begin{tcolorbox}[breakable,colback=panel,colframe=line,boxrule=0.5pt,arc=2mm,left=3mm,right=3mm,top=2mm,bottom=2mm]
\textbf{Version To Actually Say}
\medskip
\setlength{\columnsep}{1.1em}
\raggedcolumns
\begin{multicols}{2}
\RaggedRight
\setlength{\parskip}{0.7em}
\par สิ่งที่ผมพูดมาทั้งหมด จะเกิดขึ้นจริงได้อย่างไรใน 4 ปี
\par คำตอบคือ\\ เราไม่ทำทุกอย่างพร้อมกัน\\ แต่ \textbf{\textcolor{highlight}{“ติดตั้งระบบ”}} เป็นลำดับ
\par \textbf{\textcolor{highlight}{ปีที่ 1 — วางรากฐาน}}\\ สิ่งที่ต้องเกิดคือ\\ \textbf{\textcolor{highlight}{ข้อมูลชุดเดียว}} และกติกาที่ใช้ได้จริง
\par เราต้องมีระบบข้อมูลกลาง\\ ที่ทุกฝ่ายเห็นภาพเดียวกัน
\par มี แผนที่ศักยภาพ\\ รู้ว่า KU แข็งตรงไหน และควรไปต่ออย่างไร
\par และต้อง ปรับกฎและ workflow\\ ที่เป็นคอขวดให้คล่องตัว
\par ถ้าปีแรกไม่สำเร็จ\\ ปีต่อไปจะลอยทั้งหมด
\par \textbf{\textcolor{highlight}{ปีที่ 2 — ขับเคลื่อน}}\\ เมื่อเห็นภาพชัดแล้ว\\ เราจะ \textbf{\textcolor{highlight}{“เลือก”}} ไม่ใช่ \textbf{\textcolor{highlight}{“กระจาย”}}
\par ตั้ง กลุ่มเชี่ยวชาญหลัก\\ ในพื้นที่ที่มีโอกาสสร้าง impact จริง
\par วาง ระบบระดมทุนและ partner model\\ ที่สอดคล้องกับพันธกิจ
\par และเริ่ม ยกระดับบางพื้นที่สู่สากล
\par \textbf{\textcolor{highlight}{ปีที่ 3 — ขยายผล}}\\ สิ่งที่เริ่มไว้ ต้องแปลงเป็นผลเชิงสถาบัน
\par ทั้ง\\ รายได้จาก IP และกองทุน\\ ตัวชี้วัดที่สะท้อนผลลัพธ์จริง\\ และ ความน่าเชื่อถือที่คนสัมผัสได้
\par \textbf{\textcolor{highlight}{ปีที่ 4 — ฝังในระบบ}}\\ นี่คือปีที่สำคัญที่สุด
\par ไม่ใช่แค่ \textbf{\textcolor{highlight}{“อยู่ครบวาระ”}}\\ แต่คือทำให้
\par ระบบทำงานได้เอง\\ ตัวชี้วัดเชื่อมกันทุกระดับ\\ และ \textbf{\textcolor{highlight}{โปร่งใส ตรวจสอบได้จริง}}
\end{multicols}
\end{tcolorbox}

\clearpage

\section*{Slide 15: Closing}
\noindent\textcolor{muted}{Slide 15 \textbar\ Slide15.md}
\vspace{0.55em}
\begin{center}

\includegraphics[width=\linewidth,height=0.340\textheight,keepaspectratio]{/mnt/c/Users/punpi/OneDrive/KU/president/vision/qna/presentation-narrator/assets/presentation\_slides/Slide15.JPG}
\end{center}
\vspace{0.4em}
\begin{tcolorbox}[breakable,colback=panel,colframe=line,boxrule=0.5pt,arc=2mm,left=3mm,right=3mm,top=2mm,bottom=2mm]
\textbf{Version To Actually Say}
\medskip
\setlength{\columnsep}{1.1em}
\raggedcolumns
\begin{multicols}{2}
\RaggedRight
\setlength{\parskip}{0.7em}
\par ผมขอสรุปทั้งหมดในประโยคเดียวครับ
\par มหาวิทยาลัยเกษตรศาสตร์\\ ต้อง \textbf{\textcolor{highlight}{“เร่งอย่างมีระบบ”}}\\ เพื่อความยั่งยืนของคนและสังคม
\par วันนี้โลกกำลังเร่งขึ้น\\ ทั้งจาก AI การเปลี่ยนแปลงประชากร\\ การแข่งขัน และข้อจำกัดด้านทรัพยากร
\par ถ้าเรายังคิดแบบเดิม\\ ทำแบบเดิม\\ แล้วหวังผลลัพธ์ใหม่
\par ผมคิดว่าคงเป็นไปได้ยาก
\par แต่คำตอบ\\ ไม่ใช่การเร่งแบบกระจาย\\ และไม่ใช่การเปลี่ยนไปเรื่อยโดยไม่มีหลัก
\par สิ่งที่เราต้องการ\\ คือ
\par ความเร็ว\\ ที่ตั้งอยู่บน \textbf{\textcolor{highlight}{“ระบบที่ดี”}}
\par ระบบที่\\ ตัดสินใจได้อย่างมีธรรมาภิบาล\\ พัฒนาคนได้จริง\\ และสร้างความยั่งยืนให้สังคมได้
\par นี่คือเหตุผลที่ผมเสนอว่า\\ KU ต้องชัดในบทบาทของตัวเอง
\par ในฐานะ\\ \textbf{\textcolor{highlight}{The Life Systems University}}
\par และถ้าจะสรุปให้สั้นที่สุด
\par สิ่งที่เราต้องทำ คือ
\par \textbf{\textcolor{highlight}{กล้าทบทวน}}\\ \textbf{\textcolor{highlight}{จัดวางใหม่}}\\ และ\textbf{\textcolor{highlight}{ทำให้เกิดขึ้นจริง}}
\par \textbf{\textcolor{highlight}{Rethink. Reframe. Realize.}}
\end{multicols}
\end{tcolorbox}
\end{document}
